% === VIII. DISCUSSION AND LIMITATIONS ========================================
\section{Discussion and Limitations}
\label{sec:disc}

The reachable-friction bound~\eqref{eq:mureach} is kinematic---no tire
parameter, no simulator-specific quantity, no property of the residual---so it
applies to any implementation of the virtual-steady-state-data strategy
of~\cite{ontrack2025}, in simulation or on hardware. Its practical content is
a design rule, \emph{the rollout must demand at least the friction one intends
to identify}, and where it cannot, $D$ should be reported as unidentified
rather than fitted. The diagnostic signature, coefficients resting exactly on box constraints, is
likewise generic and should be logged by any bounded tire fit, being otherwise
indistinguishable from a successful one. So is the insufficiency of
per-coefficient box constraints for admissibility: both failures we
encountered---a degenerate set, and a pairwise oversteering one unstable inside
the controller's operating range, each satisfying every box constraint---are
properties of the parameterisation, not of CARLA.

What does not generalise is every number attached to this plant. The tire
values describe a PhysX model with a configured friction coefficient, so
agreement with them shows only that the identifier recovers the plant it
observes; hardware replication is necessary before any accuracy claim on a
vehicle. The gap between configured wheel
friction (\num{1.5}) and realised peak axle friction (\numrange{1.00}{1.05}),
and the \num{0.76} achieved/commanded steering ratio, are properties of this
configuration; the lessons---score against the effective friction the simulator
publishes, and identify against the achieved actuator state---are general, the
magnitudes are not.

\subsection{Limitations}
\label{subsec:limits}

\emph{One run per configuration, so nothing in Fig.~\ref{fig:comparison}
carries an uncertainty:} the \SI{0.092}{\meter} against \SI{0.177}{\meter} RMS
tracking difference has no confidence interval and no evidence that it exceeds
run-to-run variation, which is not measured; nothing prevents repetition, so
this is a gap in the evidence, not the method. \emph{Single surface, vehicle and trajectory:} friction is static at
\num{1.05}, and the surface-change adaptation experiment, arguably the
strongest argument for on-track identification, is not performed. \emph{The friction estimator's
excitation sweep is synthetic;} on the CARLA runs the warm start sees only the
excitation level the lap supplies. \emph{The two configurations differ in four
settings at once}, so Section~\ref{sec:results} attributes nothing to the
residual, the loss, the solver or the warm start individually, and with the
excitation-aware rollout active in \emph{both} arms the closed-loop comparison
is not evidence for the identifiability correction---the sweep of
Section~\ref{subsec:sweep} is. \emph{Longitudinal dynamics are out of scope,}
as in the baseline.
